% =============================================================================
% Project : THE PHILOSOPHY OF VIABILITY
% Book    : Book I — Ontology of Viability
% Book ID : VS.B01
% Part    : Part I — Identity and Manifestation
% Part ID : VS.B01.P01
% Chapter : Chapter 08 — Comparative Philosophical Examination
% Chapter ID : VS.B01.P01.C08
%
% File :
%   4_ASSETS/tables/
%   VS.B01.P01.C08.S11.SS14.P003.TL01.heidegger_claims_and_ontological_limits.tex
%
% Section : Section 11 — Martin Heidegger
%
% Asset :
%   Table TL01 — Heidegger's Claims and Ontological Limits
%
% Structural Identifier :
%   VS.B01.P01.C08.S11.SS14.P003.TL01
%
% Asset Type : TL
%
% Purpose :
%   To distinguish Heidegger's principal concepts, their
%   explanatory function, and the limits of their transfer
%   into the Ontology of Viability.
%
% Intellectual Benefit :
%   Provides a compact comparative reference and prevents
%   the conflation of Being, beings, Dasein, Identity,
%   worldhood, disclosure, and temporality.
%
% Aesthetic / Experiential Benefit :
%   Summarizes the extended discussion in a readable
%   comparative matrix.
%
% Representational Limit :
%   This table is an analytical summary. It does not claim
%   that Heidegger's work is internally uncontested.
%
% Dependencies :
%   longtable, booktabs, array, project column types.
%
% =============================================================================

\small
\renewcommand{\arraystretch}{1.18}
\setlength{\tabcolsep}{4pt}

\begin{longtable}{
    >{\raggedright\arraybackslash}p{2.15cm}
    >{\raggedright\arraybackslash}p{0.25\textwidth}
    >{\raggedright\arraybackslash}p{0.25\textwidth}
    >{\raggedright\arraybackslash}p{0.25\textwidth}
}

\caption{Heidegger's Principal Concepts and Their Ontological Limits}
\label{tab:VS.B01.P01.C08.S11.SS14.P003.TL01} \\

\toprule
\textbf{Concept}
&
\textbf{Heidegger's account}
&
\textbf{Explanatory contribution}
&
\textbf{Limit for Book I}
\\
\midrule
\endfirsthead

\toprule
\textbf{Concept}
&
\textbf{Heidegger's account}
&
\textbf{Explanatory contribution}
&
\textbf{Limit for Book I}
\\
\midrule
\endhead

\midrule
\multicolumn{4}{r}{\footnotesize\textit{Continued on next page}}\\
\endfoot

\bottomrule
\endlastfoot

Ontological difference
&
Distinction between Being and beings.
&
Prevents Being from being treated as one
more entity or highest object.
&
Not equivalent to the distinction between
Being and Identity.
\\

\addlinespace

Dasein
&
The entity that we ourselves are, whose
Being is an issue for it.
&
Clarifies the existential conditions of
the question of Being.
&
Does not provide a universal definition
of every existent.
\\

\addlinespace

Being-in-the-world
&
A unitary structure of Dasein's existence.
&
Shows the significance of worldhood,
practical involvement, and relational context.
&
Does not establish that all beings are
constituted by human practical involvement.
\\

\addlinespace

Readiness-to-hand
&
Equipment encountered in practical use
and involvement.
&
Reveals referential contexts obscured
by detached inspection.
&
A mode of encounter is not identical
with the Identity of an entity.
\\

\addlinespace

Being-with and das Man
&
Sociality and anonymous public norms
are existential structures of everydayness.
&
Shows that self-understanding is socially
and practically situated.
&
Social roles and public norms do not
by themselves determine the number
of identities.
\\

\addlinespace

Care
&
The structural whole of Dasein:
projection, thrownness, and involvement.
&
Connects facticity, possibility, and
practical engagement.
&
Not a definition of Identity or a
universal structure of all beings.
\\

\addlinespace

Authenticity
&
Taking over one's own finite possibilities.
&
Clarifies an existential mode of selfhood.
&
Not a proof of numerical identity,
substantial selfhood, or moral superiority.
\\

\addlinespace

Temporality
&
Ecstatic unity of future, having-been,
and present.
&
Shows the temporal structure of Dasein
and the horizon of its possibilities.
&
Not identical with continuity, persistence,
or identity through transformation.
\\

\addlinespace

Lichtung
&
The open region of disclosure and concealment.
&
Clarifies the relation between unconcealment
and the intelligibility of entities.
&
Not identical with Manifestation or
Ontological Light in Book I.
\\

\addlinespace

Gestell
&
The technological mode of revealing
that orders entities as standing-reserve.
&
Critiques the dominance of one mode
of disclosure.
&
Not a general rejection of technology
or a complete ontology of material beings.
\\

\addlinespace

Dwelling and place
&
Place is understood through gathering
and dwelling, not merely coordinates.
&
Challenges a purely geometrical account
of spatial situatedness.
&
Gathering does not by itself establish
the Identity of a place.
\\

\end{longtable}

\normalsize

% =============================================================================
% END OF TABLE TL01
% =============================================================================